\documentclass[11pt]{article}
\usepackage[a4paper,margin=1in]{geometry}
\usepackage{amsmath,amssymb}
\usepackage{booktabs}
\usepackage{graphicx}
\usepackage{hyperref}
\hypersetup{colorlinks=true,linkcolor=blue,citecolor=blue,urlcolor=blue}
\usepackage[numbers,sort&compress]{natbib}

\newcommand{\doi}[1]{DOI: \href{https://doi.org/#1}{#1}}

\title{The Cost of Reliability in Optimal Magnetization Switching:\\
A Longitudinal-Field Front for Van der Waals Magnets}
\author{Felipe Santib\'a\~nez-Leal\\
\small CAOS open-research programme, Santiago, Chile}
\date{2026}

\begin{document}
\maketitle

\begin{abstract}
Optimal control theory delivers magnetic-field pulses that reverse a magnetic moment for the least
dissipated energy, and recent work shows these pulses reach picosecond, energy-efficient switching in
two-dimensional van der Waals magnets. Because the optimal pulse is always perpendicular to the moment,
it does nothing to damp the perturbations that thermal fluctuations excite, and the switching becomes
unreliable exactly where it matters. A longitudinal field parallel to the moment removes that
instability and is invisible to the leading-order dynamics, so it is often described as free. It is not
free: it costs the square of its amplitude integrated over the pulse. We compute, for the first time,
the cost-reliability front of that longitudinal field. Using an open-source engine that solves optimal
control over the Landau-Lifshitz-Gilbert equation and a thermostat validated against the Boltzmann
distribution, we reproduce the published stabilization physics for CrSBr, the half-hyperbolic bare
path and the counterintuitive success dip near a half anisotropy field, and we add the missing axis:
the added switching cost, which grows quadratically and reaches roughly $1.5\times10^{-10}$~T$^2$~s at
twice the anisotropy field for a demanding thermal-stability factor of twenty. The front makes the
reliability-versus-energy trade explicit and quantitative, and it is reproducible from committed
artifacts and a live instance.
\end{abstract}

\section{Introduction}

The energy cost of writing a magnetic bit has become a design constraint as data-centric computing
grows. Conventional field-driven reversal is slow and expensive, which is why current-driven
spin-transfer and spin-orbit torque dominate modern magnetic memory. Optimal control theory (OCT)
reopens the field-driven route: by shaping the applied field in time one can reverse the moment in the
picosecond range at energies up to two orders of magnitude below conventional field protocols, and in
two-dimensional van der Waals magnets this becomes competitive with current-based
approaches~\citep{badarneh2026}. The formalism rests on inverting the equation of motion so the applied
field is a function of the trajectory, which turns a constrained optimization into an unconstrained one
over the reversal path~\citep{kwiatkowski2021}.

A feature of the optimal pulse is that it is always perpendicular to the
moment~\citep{kwiatkowski2021}. That perpendicularity is efficient, but it means the pulse exerts no
restoring torque on the transverse perturbations that thermal noise continuously excites. Linearizing
the dynamics about the optimal control path shows that the perturbation growth is governed by the
eigenvalues of the energy Hessian, and that a large part of a bare reversal is dynamically unstable
(hyperbolic), which is the primary cause of the pulse and the moment losing phase
lock~\citep{badarneh2023stability}. A constant longitudinal field, parallel to the instantaneous moment
direction, shifts those eigenvalues and removes the hyperbolic domain, driving the switching success
rate to unity.

Because the optimal pulse is perpendicular, the longitudinal field does no work at leading order and is
frequently treated as a cost-free stabilizer. This is misleading. The switching cost is the time
integral of the squared total field, so a longitudinal component of amplitude $B_r$ adds exactly
$\int_0^T B_r^2\,dt$ to the cost. The trade between reliability and energy is therefore real, and to
our knowledge it has not been quantified. We quantify it here.

\section{Model and methods}

\subsection{Dynamics and cost}

We consider a single-domain moment with easy-axis anisotropy energy $K$, magnetic moment $\mu$,
gyromagnetic ratio $\gamma$ and Gilbert damping $\alpha$. The Landau-Lifshitz-Gilbert (LLG) equation in
Gilbert form is
\begin{equation}
(1+\alpha^2)\,\dot{\vec s} = -\gamma\,\vec s \times (\vec b_i + \vec b)
-\alpha\gamma\,\vec s \times \bigl[\vec s \times (\vec b_i + \vec b)\bigr],
\end{equation}
with $\vec b_i = -\mu^{-1}\partial E/\partial \vec s$ the internal field and $\vec b$ the applied
field. The switching cost, proportional to the Joule heating of the circuit that generates the field,
is
\begin{equation}
\Phi = \int_0^T |\vec b(t)|^2\, dt,
\label{eq:cost}
\end{equation}
in units of T$^2$~s. Equation~\eqref{eq:cost} is not an energy: converting it to joules requires a
circuit constant with units of joules per T$^2$~s, which we do not assume here, so all costs are
reported in their model-free units.

\subsection{The optimal pulse and its stability}

For a uniaxial magnet, $E = -K s_z^2$, the optimal control path and its perpendicular pulse are known
in closed form as Jacobi elliptic functions~\citep{kwiatkowski2021}. Linearizing the LLG about that
path gives a two-component transverse perturbation whose growth is set by
\begin{equation}
w_1 = B_r + \frac{K}{\mu}\cos 2\theta, \qquad
w_2 = B_r + \frac{K}{\mu}\cos^2\theta,
\end{equation}
where $\theta$ is the polar angle along the path and $B_r$ is the applied-field component parallel to
the moment~\citep{badarneh2023stability}. Perturbations are bounded (elliptic) when $w_1 w_2 > 0$ and
divergent (hyperbolic) when $w_1 w_2 \le 0$. At $B_r = 0$ the reversal is hyperbolic for
$\pi/4 < \theta < 3\pi/4$, exactly half the path. Adding $|B_r| > K/\mu$ removes the hyperbolic domain.

\subsection{Thermostat and success rate}

We integrate an ensemble of stochastic LLG trajectories with a fluctuating thermal field whose
strength is fixed by the fluctuation-dissipation theorem, each Cartesian component drawn per step from
a Gaussian of variance $2\alpha k_B T/(\gamma\mu\,dt)$~\citep{evans2014}, using a norm-preserving
stochastic Heun scheme. The thermostat is validated against the Boltzmann distribution it must
reproduce: near a uniaxial minimum the mean square polar angle is $\langle\theta^2\rangle = k_B T/K$,
and the thermostat recovers it to within ten percent. For each longitudinal field $B_r$ we apply the
bare optimal pulse together with a constant $B_r$ along the instantaneous moment direction, run the
ensemble at temperature, and record the fraction that ends in the reversed basin (the success rate) and
the added cost $B_r^2 T$.

\subsection{Material and engine}

The material is monolayer CrSBr, whose measured exchange Hamiltonian and triaxial, substrate-tunable
anisotropy make it the natural host of the internal-torque mechanisms of this
family~\citep{scheie2022,rudenko2023}. Its magnetic parameters are curated from the primary sources
with provenance. All computation is done by \texttt{spinoct}, an open-source (MIT) engine that solves
optimal control over LLG dynamics; the results here are reproducible from committed artifacts.

\section{Results}

We compute the cost-reliability front for CrSBr at a thermal stability factor
$\Delta = K/(k_B T) = 20$, a demanding, sub-memory-grade regime chosen so the reliability gain is
visible. Table~\ref{tab:front} reports the hyperbolic fraction of the path, the Monte-Carlo success
rate over $400$ copies with its binomial confidence interval, and the added cost, as the longitudinal
field is swept in units of the anisotropy field $K/\mu$.

\begin{table}[h]
\centering
\caption{The longitudinal-field cost-reliability front for CrSBr at $\Delta = 20$. The bare path
($B_r = 0$) is half-hyperbolic; a field of order the anisotropy field removes the instability; the
added cost grows as $B_r^2$. Success rates are over $400$ stochastic copies.}
\label{tab:front}
\begin{tabular}{cccc}
\toprule
$B_r/(K/\mu)$ & hyperbolic fraction & success rate & added cost (T$^2$~s) \\
\midrule
0.0 & 0.50 & $1.000 \pm 0.005$ & $0.00$ \\
0.5 & 0.33 & $0.980 \pm 0.014$ & $6.1\times10^{-12}$ \\
1.0 & 0.00 & $0.958 \pm 0.020$ & $2.5\times10^{-11}$ \\
1.5 & 0.00 & $0.998 \pm 0.005$ & $5.5\times10^{-11}$ \\
2.0 & 0.00 & $1.000 \pm 0.005$ & $9.8\times10^{-11}$ \\
2.5 & 0.00 & $1.000 \pm 0.005$ & $1.5\times10^{-10}$ \\
\bottomrule
\end{tabular}
\end{table}

Three features are worth drawing out. First, the hyperbolic fraction falls from one half at $B_r = 0$
to zero once $B_r \ge K/\mu$, exactly as the linearized stability analysis
predicts~\citep{badarneh2023stability}. Second, the success rate shows a shallow, counterintuitive dip
around $B_r \approx K/\mu$ before recovering to unity, the same non-monotonicity reported in the
stabilization study, which arises because the eigenvalue ratio passes through $-1$ near the barrier top
in that range. Third, and new here, the added cost grows quadratically with the field, reaching
$1.5\times10^{-10}$~T$^2$~s at $B_r = 2.5\,K/\mu$. This is the price of driving the reliability from a
partly unstable bare pulse to certain switching.

The consequence for device design is that the longitudinal field is not a free stabilizer to be turned
up without limit. The front in Table~\ref{tab:front} lets a designer read off the smallest field that
reaches a target reliability and pay only that cost. For CrSBr at $\Delta = 20$, a field of about one
anisotropy field already removes the dynamical instability, and the residual success-rate deficit is
closed by a modest further increase at a cost that remains well below the bare switching cost.

\section{Discussion}

The longitudinal-field front reframes reliability as a purchasable quantity with a stated price, rather
than a binary the optimal pulse either has or lacks. The bare optimal pulse minimizes energy but is
dynamically fragile; the stabilized pulse is reliable but costs more; the front is the locus of the
efficient trades between them. Because the added cost is separable, $\int B_r^2\,dt$, and the bare cost
is the closed-form optimal-control cost, the whole front is cheap to compute once the stability
eigenvalues, which come from the same Hessian the optimal-control solver already uses, are in hand.

This suggests a deterministic surrogate for the expensive stochastic sweep. The hyperbolicity integral
$\int_0^T \max(0, -w_1 w_2)\,dt$ along a path is a purely deterministic measure of dynamical fragility
that requires no ensemble. Whether penalizing it predicts the Monte-Carlo success rate, and therefore
whether a single deterministic optimal-control solve can target a reliability without any stochastic
simulation, is a well-posed question this front makes answerable and which we leave to future work.

\section{Conclusion}

We have quantified the cost of reliability in optimal magnetization switching. For monolayer CrSBr, a
longitudinal field of order the anisotropy field removes the dynamical instability of the bare optimal
pulse and drives the switching success rate to unity, at an added switching cost that grows
quadratically with the field and that we report as a front. The reliability-versus-energy trade is now
explicit and quantitative, computed by an open engine and reproducible from committed artifacts and a
live instance at \url{https://espira.fasl-work.com}.

\section*{Data and code availability}

The engine \texttt{spinoct} is open source (MIT) at
\url{https://github.com/fsantibanezleal/CAOS_SpinOCT}. The product, the curated material parameter
database, the canonical bake and the companion web application are open (MIT) at
\url{https://github.com/fsantibanezleal/CAOS_RES_Espira}, with a live instance at
\url{https://espira.fasl-work.com}. Every number in Table~\ref{tab:front} is reproduced from a
committed artifact.

\begin{thebibliography}{9}
\bibitem{badarneh2026} M. H. Badarneh, P. Cai, E. J. G. Santos,
\emph{Optimal Control Drives Ultrafast and Energy-Efficient Magnetization Switching in Van der Waals
Magnets}, Advanced Materials \textbf{e23059} (2026). \doi{10.1002/adma.202523059}.
\bibitem{kwiatkowski2021} G. J. Kwiatkowski, M. H. A. Badarneh, D. V. Berkov, P. F. Bessarab,
\emph{Optimal Control of Magnetization Reversal in a Monodomain Particle by Means of Applied Magnetic
Field}, Phys. Rev. Lett. \textbf{126}, 177206 (2021). \doi{10.1103/PhysRevLett.126.177206}.
\bibitem{badarneh2023stability} M. H. A. Badarneh, G. J. Kwiatkowski, P. F. Bessarab,
\emph{Enhancing thermal stability of optimal magnetization reversal in nanoparticles},
arXiv:2312.11293 (2023). \doi{10.48550/arXiv.2312.11293}.
\bibitem{evans2014} R. F. L. Evans, W. J. Fan, P. Chureemart, T. A. Ostler, M. O. A. Ellis,
R. W. Chantrell, \emph{Atomistic spin model simulations of magnetic nanomaterials}, J. Phys.: Condens.
Matter \textbf{26}, 103202 (2014). \doi{10.1088/0953-8984/26/10/103202}.
\bibitem{scheie2022} A. Scheie, M. Ziebel, D. G. Chica, et al.,
\emph{Spin Waves and Magnetic Exchange Hamiltonian in CrSBr}, Advanced Science \textbf{9}, 2202467
(2022). \doi{10.1002/advs.202202467}.
\bibitem{rudenko2023} A. N. Rudenko, M. R\"osner, M. I. Katsnelson,
\emph{Dielectric tunability of magnetic properties in orthorhombic ferromagnetic monolayer CrSBr},
arXiv:2302.12672 (2023). \doi{10.48550/arXiv.2302.12672}.
\end{thebibliography}

\end{document}
